%!TeX root=../Memory.tex

\chapter{Desenvolupament}
\label{cap:desenvolupament}

Aquest capítol descriu com s'ha executat el projecte durant la seva fase de
construcció. Mentre que el capítol~\ref{cap:planificacio} ha exposat la
\emph{planificació} (cicle de vida, jerarquia del \emph{backlog}, \aclu{DoR}
i \aclu{DoD}, política de versionat), aquí es relata el desenvolupament real:
com s'ha organitzat l'equip dia a dia, com s'ha materialitzat la jerarquia
del \emph{backlog} en \emph{issues} concretes, què ha entregat cada
\emph{sprint}, quines versions intermèdies s'han generat, quins problemes
s'han trobat i com s'han demostrat les funcionalitats. L'objectiu és deixar
constància fidel del \emph{què} ha passat, no del \emph{què s'havia
previst}.

%-----------------------------------------------------------------------
\section{Organització de l'equip}
\label{sec:organitzacio-equip}

L'equip està format per quatre estudiants del Grau d'Enginyeria
Informàtica de la \ac{UIB}, sense experiència prèvia en
\emph{blockchain} ni en criptografia aplicada (vegeu
secció~\ref{sec:equip} del capítol~\ref{cap:introduccio}). Aquesta
homogeneïtat de perfils ha facilitat la presa de decisions
col·lectiva i la revisió creuada del codi, però ha obligat l'equip a
adquirir des de zero el coneixement tècnic necessari (vegeu
secció~\ref{sec:problemes}).

\subsection{Rols i responsabilitats}

Cada membre de l'equip ha liderat una àrea funcional del producte
durant tot el projecte, mantenint la mateixa responsabilitat des de
la fase de recerca fins a l'entrega final. Aquesta estabilitat ha
permès acumular coneixement profund sobre cada paquet de treball i
evitar el cost de canvis freqüents de context. La
taula~\ref{tab:rols-desenv} sintetitza els rols.

\begin{table}[ht]
\centering
\small
\begin{tabular}{@{}p{3.3cm}p{4cm}p{5cm}@{}}
\toprule
\textbf{Membre} & \textbf{Àrea principal} & \textbf{Responsabilitats concretes} \\
\midrule
Jordi Vanyó Mestre & Client web i integració Algorand & Interfície React, integració amb \texttt{algosdk}, codificació i descodificació ARC-4, lectura de \emph{boxes}, signatura de transaccions des del navegador. \\
\addlinespace
Dylan Luigi Canning Garcia & Ancoratge i \emph{NotaryContract} & Servei Python d'ancoratge (\texttt{network/anchoring/}), \emph{hash} determinista (JSON canònic + \ac{SHA}-256), submissió al contracte d'Ethereum, contracte \texttt{NotaryContract} a Sepolia. \\
\addlinespace
Antonio Contestí Coll & Coordinació i documentació & Gestió del \emph{backlog} a GitHub Projects, seguiment de \emph{sprints}, coherència de la documentació tècnica, redacció de la memòria. \\
\addlinespace
Marc Link Cladera & \emph{Smart contract} Algorand i arquitectura & Disseny i implementació del contracte \texttt{Demochain} en \emph{algopy}, diagrames C4 \cite{brown2018c4}, redacció dels \acp{ADR}. \\
\bottomrule
\end{tabular}
\caption[Rols i responsabilitats]{Distribució de rols i responsabilitats
durant el desenvolupament. Cap rol ha estat exclusiu: totes les \acsp{PR}
han passat per revisió creuada (vegeu
secció~\ref{sec:versionat} del capítol~\ref{cap:planificacio}).}
\label{tab:rols-desenv}
\end{table}

Aquesta distribució per àrees no ha estat estricta. La regla de
protecció de \texttt{main} --que impedeix \emph{mergejar} la pròpia
\ac{PR}-- ha forçat que cada membre revisi codi de les altres àrees i,
en conseqüència, hi adquireixi un coneixement funcional bàsic. Les
decisions arquitectòniques rellevants s'han pres en reunions setmanals
conjuntes, inicialment presencials i progressivament híbrides.

\subsection{Dinàmica de treball}

L'equip ha funcionat amb una cadència setmanal estable. Cada setmana
s'ha celebrat una reunió de coordinació amb tres punts fixos: revisió
del que s'ha fet, identificació de bloqueigs i planificació de la
setmana següent. A més, dins de cada \emph{sprint} s'ha celebrat una
\emph{retrospectiva} curta per identificar millores de procés. Les
decisions de calat tècnic s'han documentat als \acp{ADR} corresponents
(\ac{ADR} 000 sobre \emph{algopy} i \ac{ADR} 001 sobre la plataforma
\emph{blockchain}), de manera que la discussió no s'ha hagut de
repetir.

%-----------------------------------------------------------------------
\section{Organització del backlog}
\label{sec:backlog-desenv}

El \emph{backlog} del projecte s'ha gestionat amb \textbf{GitHub
Projects} sobre el repositori \texttt{jvanyom/demochain}
\cite{demochainrepo}, seguint l'estructura jeràrquica de cinc nivells
descrita a la secció~\ref{sec:jerarquia} del
capítol~\ref{cap:planificacio}. Aquesta secció n'il·lustra l'aplicació
real amb exemples concrets del repositori.

\subsection{Materialització de la jerarquia}
\label{sec:materialitzacio-jerarquia}

La cadena \textbf{èpica $\rightarrow$ història d'usuari / \emph{spike}
$\rightarrow$ tasca} s'ha aplicat de manera consistent durant tot el
projecte. La figura conceptual següent mostra un exemple complet
extret del \emph{backlog}:

\begin{itemize}
    \item \textbf{Èpica E2 -- Propostes i aprovació} (\emph{issue} \#5).
    Conjunt de funcionalitats per a la creació, validació i aprovació
    de propostes per quòrum de dos terços.
    \begin{itemize}
        \item \textbf{Història d'usuari E2-US1} -- \emph{Com a membre
        del cens vull crear una proposta dins de la meva organització}.
        Inclou els criteris d'acceptació en format
        \emph{Donat/Quan/Aleshores} sobre validacions de dates,
        opcions mínimes i pertinença al cens.
        \begin{itemize}
            \item \textbf{Tasca E2-US1-T3} -- Implementació del mètode
            \texttt{create\_proposal} al contracte amb validació de
            la finestra temporal mínima (3 dies abans de l'inici, 1
            dia mínim de votació) i del nombre d'opcions ($\geq 2$).
        \end{itemize}
    \end{itemize}

    \item \textbf{\emph{Spike} E1-SP1} -- \emph{Quin llenguatge cal
    triar per al contracte d'Algorand: TEAL, PyTEAL o algopy?} El
    resultat d'aquest \emph{spike} s'ha materialitzat al document
    \ac{ADR} 000 (\texttt{docs/decisions/000-smart-contract-language.md}).
\end{itemize}

A cada \emph{issue} se li ha assignat, a més del tipus
(\texttt{type/epic}, \texttt{type/user-story}, \texttt{type/spike},
\texttt{type/task}, \texttt{type/subtask}), una etiqueta d'àrea
funcional: \texttt{area/smart-contract}, \texttt{area/frontend},
\texttt{area/network}, \texttt{area/census}, \texttt{area/wallet},
\texttt{area/test}, \texttt{area/research}, \texttt{area/docs} i
\texttt{area/devops}. Aquesta combinació de tipus i àrea ha permès
filtrar el \emph{backlog} segons diferents perspectives durant les
reunions de planificació.

\subsection{Volum i estat final del backlog}

A l'entrega final, el \emph{backlog} comprèn \textbf{nou èpiques} i més
de \textbf{cent \emph{issues}} traçades, tal com es detalla a la
taula~\ref{tab:epiques} del capítol~\ref{cap:planificacio}. La majoria
de les èpiques estan tancades; només dues queden obertes
explícitament:

\begin{itemize}
    \item \textbf{E5 -- Elegibilitat de votants amb \ac{ZK}}
    (\emph{issue} \#340), que depèn íntegrament de la capa
    criptogràfica no implementada.
    \item \textbf{E11 -- Setup del projecte i del repositori}
    (\emph{issue} \#402), oberta com a \emph{epic} contenidor de
    tasques de manteniment continu del repositori.
\end{itemize}

Els tres últims \emph{sprints} es registren formalment com a
\emph{milestones} a GitHub, amb les \emph{issues} corresponents
assignades a cadascun. Una de les retrospectives ha identificat que
els tres primers \emph{sprints} no es van \emph{milestonar} de
manera estricta i caldria assignar-hi els \emph{issues} corresponents
per a la traçabilitat completa. Aquesta tasca queda recollida com a
punt pendent a l'\emph{epic} E11.

%-----------------------------------------------------------------------
\section{Resum de les iteracions de desenvolupament}
\label{sec:iteracions}

El desenvolupament s'ha estructurat en sis iteracions
(\emph{sprints}) de dues setmanes. Aquesta secció en relata
l'execució real, complementant la planificació exposada a la
secció~\ref{sec:sprints} del capítol~\ref{cap:planificacio}.

\subsection{Sprint 1 -- Recerca preliminar (febrer-març)}

Iteració íntegrament dedicada a la fase de recerca. L'equip ha
estudiat el funcionament d'Algorand
\cite{chen2017algorand,algorandfoundation2024algokit}, les
alternatives a Ethereum \cite{buterin2014ethereum} i els mètodes de
votació ordinal. S'ha aixecat un primer \emph{LocalNet} d'Algorand i
s'ha provat el desplegament d'un contracte mínim per familiaritzar-se
amb la cadena d'eines. \textbf{Entrega:} primera versió del repositori
amb prova de concepte sobre LocalNet i tres documents de recerca
(\texttt{docs/research/algorand.md}, \texttt{ethereum.md},
\texttt{voting-methods-comparison.md}).

\subsection{Sprint 2 -- Decisions arquitectòniques (març)}

S'han fixat les decisions tècniques clau del projecte: tria del
llenguatge del contracte (\emph{algopy} amb compilador Puya
\cite{algorandfoundation2024puya}, formalitzada a l'\ac{ADR} 000) i
tria de la plataforma \emph{blockchain} (Algorand permissionada amb
ancoratge a Ethereum, formalitzada a l'\ac{ADR} 001). Paral·lelament,
s'han redactat els primers diagrames C4 \cite{brown2018c4} i s'ha
seleccionat el mètode de Schulze \cite{schulze2011newmethod} com a
algorisme de recompte, amb el corresponent document de recerca
(\texttt{docs/research/schulze.md}). \textbf{Entrega:} primer esborrany
del contracte \emph{algopy} amb les operacions bàsiques d'organització
i diagrames C4 a nivells 1 i 2.

\subsection{Sprint 3 -- Primer prototip de votació (març -- 12 abril)}

Implementació del primer prototip funcional. Al contracte d'Algorand
s'hi ha afegit la lògica de propostes (creació, validació de dates i
opcions) i una versió inicial de votació per pluralitat. Al client
web s'ha construït el \emph{mockup} de les pantalles principals amb
React i \emph{TailwindCSS}, i s'ha configurat el primer
\texttt{docker-compose.yml} amb \texttt{algod}, indexer i client.
\textbf{Entrega intermèdia} a la convocatòria del 12 d'abril.

Aquest \emph{sprint} acaba amb la \emph{sessió d'estimació realista}
en què l'equip pren la decisió de \textbf{reduir l'abast del \ac{MVP}}
i deixar la criptografia avançada com a treball futur (vegeu
secció~\ref{sec:problemes} i secció~\ref{sec:reduccio-abast} del
capítol~\ref{cap:planificacio}).

\subsection{Sprint 4 -- Schulze, resultats i integració (17--30 abril)}

Aquest \emph{sprint} consolida la presència del repositori actual i
incorpora la majoria del codi del client. S'implementa el mòdul de
Schulze en TypeScript sobre la matriu de comparació per parelles
(\emph{issues} \#132 i \#133) amb l'algorisme de Floyd-Warshall
\cite{floyd1962algorithm}, la pàgina de resultats amb guanyador i
\emph{rànquing} (\#136), les matrius de comparació per parelles
col·lapsables (\#139), la targeta \emph{«El meu vot»} al detall de la
proposta (\#144) i el panell de verificació pública del vot (\#145).
\textbf{Entrega:} versió integrada client + contracte capaç d'executar
el cicle complet d'una elecció sobre LocalNet.

\subsection{Sprint 5 -- Testing, refinament i robustesa (1--14 maig)}

\emph{Sprint} orientat a la qualitat i estabilitat. S'incorporen
cobertures de proves unitàries del domini de propostes (\#458), tests
dels \emph{mappers} de consultes Algorand (\#457) i tests d'integració
del \emph{governance client} contra LocalNet (\#456). Es corregeix
el \emph{bug} d'accés a transaccions via Lora (\#463) i es refina la
pantalla de votació amb edició d'hora i minuts mitjançant teclat
(\#470). \textbf{Entrega:} cinc fitxers de test unitari al contracte
(\texttt{voting-contract/tests/}: \texttt{approval\_unit\_test.py},
\texttt{cast\_election\_vote\_unit\_test.py},
\texttt{census\_unit\_test.py}, \texttt{organization\_unit\_test.py},
\texttt{proposal\_unit\_test.py}) i tests al client.

\subsection{Sprint 6 -- Documentació, demo i entrega (15--25 maig)}

\emph{Sprint} final dedicat a redacció de la memòria, preparació de la
presentació oral, retrospectiva i empaquetament. S'han generat els
diagrames C4 a nivell 3 (un per contenidor) i els diagrames de
seqüència Mermaid dels fluxos principals. \textbf{Entrega:} versió
\texttt{v1.0.0} del producte i memòria final.

%-----------------------------------------------------------------------
\section{Versions generades i entregades}
\label{sec:versions}

A banda de la versió final del producte, durant tot el projecte s'han
generat entregables intermedis amb cadència setmanal. Aquesta secció
n'inventaria els més rellevants, distingint entre entregables de
recerca, entregables d'enginyeria i la versió final del producte.

\subsection{Recerca obtinguda setmanalment}
\label{sec:recerca-setmanal}

Durant tot el projecte, l'equip ha mantingut una pràctica setmanal de
documentar el coneixement adquirit. Aquesta pràctica ha estat
especialment intensa durant la fase de recerca (\emph{sprints} 1 i 2),
però s'ha mantingut tot al llarg del projecte cada vegada que ha
sorgit un \emph{spike} de recerca. Els resultats es materialitzen en
dos llocs del repositori:

\begin{itemize}
    \item \textbf{\texttt{docs/research/}}: documents tècnics amb les
    conclusions d'estudis comparatius. A l'entrega final n'hi ha
    quatre:
    \begin{itemize}
        \item \texttt{algorand.md} -- Estudi en profunditat
        d'Algorand: model de consens, \ac{AVM}, \emph{boxes}, ARC-4
        \cite{algorandfoundation2023arc4} i ARC-56
        \cite{algorandfoundation2023arc56}.
        \item \texttt{ethereum.md} -- Anàlisi d'Ethereum com a opció
        alternativa i com a destinació de l'ancoratge, incloent-hi
        \ac{EIP}-1559 \cite{eip1559}.
        \item \texttt{voting-methods-comparison.md} -- Comparativa de
        mètodes de votació ordinal (Borda, Condorcet, IRV, Schulze)
        amb les propietats formals respectives.
        \item \texttt{schulze.md} -- Estudi en detall del mètode de
        Schulze \cite{schulze2011newmethod} i del seu càlcul amb
        Floyd-Warshall \cite{floyd1962algorithm}.
    \end{itemize}

    \item \textbf{\texttt{docs/decisions/}}: \aclp{ADR} amb les
    decisions arquitectòniques formals derivades dels estudis
    anteriors:
    \begin{itemize}
        \item \ac{ADR} 000 -- Tria del llenguatge de \emph{smart
        contract} (\texttt{000-smart-contract-language.md}).
        \item \ac{ADR} 001 -- Tria de la plataforma \emph{blockchain}
        (\texttt{001-blockchain-choice.md}).
    \end{itemize}
\end{itemize}

Aquesta pràctica ha resultat especialment valuosa: cada document de
recerca i cada \ac{ADR} ha servit per evitar que la mateixa discussió
es repetís en \emph{sprints} posteriors i ha proporcionat justificació
formal per a les decisions de disseny.

\subsection{Versions intermèdies del producte}

Pel que fa al producte de programari, s'han generat dues versions
significatives intermèdies abans del lliurament final:

\begin{itemize}
    \item \textbf{Prototip Sprint 3} (12 d'abril): \emph{mockup} del
    client i contracte amb votació per pluralitat. Entrega intermèdia
    formal a la convocatòria del 12 d'abril.
    \item \textbf{Integració Sprint 4} (30 d'abril): primera versió
    capaç d'executar el cicle complet d'una elecció (organització,
    proposta, aprovació, votació preferencial, recompte Schulze) sobre
    LocalNet.
\end{itemize}

\subsection{Versió final entregada}

La versió final entregada és la \textbf{\texttt{v1.0.0}}, generada al
final del Sprint 6. Comprèn:

\begin{itemize}
    \item \emph{Smart contract} d'Algorand \texttt{Demochain} en
    \emph{algopy} (\texttt{voting-contract/}).
    \item \emph{Smart contract} d'Ethereum \texttt{NotaryContract} en
    Solidity (\texttt{network/ethereum/}).
    \item Servei d'ancoratge en Python (\texttt{network/anchoring/}).
    \item Client web React 19 + TypeScript (\texttt{client/}).
    \item Infraestructura Docker Compose (\texttt{docker-compose.yml},
    \texttt{conduit.yml}).
    \item Documentació completa: \texttt{docs/} amb subdirectoris
    \texttt{research/}, \texttt{decisions/}, \texttt{diagrams/} i
    \texttt{memory/}.
\end{itemize}

%-----------------------------------------------------------------------
\section{Problemes trobats i la seva resolució}
\label{sec:problemes}

Durant el desenvolupament s'han trobat diversos obstacles. Aquesta
secció recull els dos més rellevants pel seu impacte sobre l'abast i
sobre la planificació, ja que la majoria d'incidències tècniques
puntuals queden documentades a les \emph{issues} de GitHub i als
\emph{pull requests} corresponents.

\subsection{Criptografia avançada fora de l'abast del MVP}
\label{sec:criptografia-fora}

El problema de planificació més important del projecte ha estat la
constatació que la \textbf{capa de criptografia avançada} prevista
inicialment no era assolible dins l'horitzó temporal del \ac{TFG}.
L'arquitectura inicial preveia:

\begin{itemize}
    \item Vots xifrats amb ElGamal exponencial homomòrfic.
    \item Proves \acsp{zkSNARK} Groth16
    \cite{cortier2019belenios,adida2008helios} per acreditar
    l'elegibilitat del votant sense revelar-ne la identitat.
    \item Cerimònia de generació distribuïda de claus (DKG) i
    desxifrat amb \emph{trustees} amb proves Chaum-Pedersen.
    \item Integració amb un \ac{IdP} extern (Keycloak) per a
    l'autenticació prèvia del votant.
\end{itemize}

\paragraph{Detecció del problema.} La detecció es va produir al final
del Sprint~3, en una sessió d'estimació realista en què es va
quantificar el cost d'implementar cada primitiva. Les estimacions
indicaven que sense un mes addicional de recerca específica per
primitiva, l'equip no tenia les bases per construir-les amb les
garanties necessàries (vegeu
secció~\ref{sec:reduccio-abast} del
capítol~\ref{cap:planificacio}).

\paragraph{Resolució adoptada.} La decisió, presa en grup i
documentada a les \emph{issues} corresponents, va ser:

\begin{enumerate}
    \item Reduir el \ac{MVP} a votació \textbf{en clar} amb
    verificació pública via consulta directa a la cadena Algorand i
    ancoratge K-de-N a Ethereum.
    \item Tancar les èpiques E3, E4 i E6 amb el subconjunt de
    funcionalitats efectivament implementades (votació preferencial,
    recompte Schulze, connexió de \emph{wallet}) i deixar la resta de
    \acsp{US} obertes com a treball futur explícit.
    \item Mantenir oberta l'èpica E5 sencera, ja que tota la seva
    funcionalitat depèn de la capa criptogràfica.
\end{enumerate}

\paragraph{Lliçó apresa.} Vista en retrospectiva, la decisió és
coherent amb les recomanacions de la literatura: davant un objectiu
ambiciós i un horitzó tancat, és preferible entregar un subconjunt
funcional i sòlid que un superconjunt incomplet o inestable. La
criptografia avançada queda identificada explícitament al capítol de
conclusions com a línia de treball futura.

\subsection{Falta de coneixement inicial}
\label{sec:falta-coneixement}

Cap membre de l'equip tenia experiència prèvia en \emph{blockchain}
ni en criptografia aplicada. Aquesta absència de base ha estat el
segon problema rellevant del projecte: cap línia de codi del
contracte podia escriure's abans d'entendre el model d'execució de
l'\ac{AVM}, l'estructura de \emph{boxes} d'emmagatzematge, el sistema
de tipus d'\emph{algopy} i les convencions ARC-4
\cite{algorandfoundation2023arc4}.

\paragraph{Resolució adoptada.} Es va adoptar una \textbf{fase de
recerca i fonaments explícita} (febrer-abril de 2026), durant la qual
no es va comprometre cap entrega de producte. La fase va consistir en:

\begin{itemize}
    \item Lectura sistemàtica de la documentació oficial d'Algorand,
    AlgoKit \cite{algorandfoundation2024algokit} i Puya
    \cite{algorandfoundation2024puya}.
    \item Redacció dels documents de recerca a \texttt{docs/research/}
    com a mecanisme de consolidació del coneixement.
    \item Formalització de les decisions clau als \acp{ADR} 000 i
    001, perquè el coneixement adquirit quedés codificat i no s'hagués
    de tornar a discutir.
    \item Provatura pràctica: aixecament d'un \emph{LocalNet}
    d'Algorand i desplegament d'un contracte mínim per familiaritzar-se
    amb la cadena d'eines.
\end{itemize}

\paragraph{Resultat.} El cost d'oportunitat d'aquesta fase ha estat
alt --un mes sense codi de producte entregat-- però en retrospectiva
ha resultat la decisió més encertada del projecte: el contracte que
es va escriure després estava lliure d'errors conceptuals que
altrament haurien requerit reescriptures costoses.

%-----------------------------------------------------------------------
\section{Demostracions funcionals}
\label{sec:demos}

Cada \emph{sprint} ha culminat amb una demostració funcional interna
de les funcionalitats incorporades. Aquesta secció en sintetitza
l'evolució i descriu la demostració de la versió final.

\subsection{Demostracions intermèdies}

\begin{itemize}
    \item \textbf{Sprint 3 (12 abril).} Demostració del prototip de
    votació per pluralitat: creació d'organització, creació de
    proposta, votació simple i visualització del resultat sobre el
    \emph{mockup} del client.

    \item \textbf{Sprint 4 (30 abril).} Demostració del cicle complet
    sobre LocalNet: creació d'organització amb cens, redacció d'una
    proposta, fase d'aprovació per quòrum, votació preferencial amb
    \emph{rànquing} arrossegable, càlcul del guanyador amb Schulze i
    visualització de la matriu de comparació per parelles.

    \item \textbf{Sprint 5 (14 maig).} Demostració de la robustesa:
    execució de tota la bateria de tests del contracte, càrrega
    massiva de cens via \ac{CSV} amb trossejament automàtic en lots
    de set adreces, panell de verificació pública del vot per a una
    adreça arbitrària, i correcció del \emph{bug} d'accés a
    transaccions via Lora.
\end{itemize}

\subsection{Demostració de la versió final}
\label{sec:demo-final}

La demostració de la versió \texttt{v1.0.0} a l'entrega final cobreix
el cicle complet d'una elecció sobre l'entorn local i és
\textbf{reproductible} per qualsevol persona amb el repositori i
Docker.

\paragraph{Preparació de l'entorn.} Una sola comanda aixeca tota la
infraestructura: \texttt{docker compose up}. Això arrenca
seqüencialment el node Algorand (\emph{lead} i \emph{follower}), el
desplegament automàtic del contracte (que persisteix el \texttt{APP\_ID}
a \texttt{client/.env.local}), la base de dades PostgreSQL, Conduit,
l'\emph{indexer} i el client web servit per Vite (vegeu
secció~\ref{sec:cicle-vida} del capítol~\ref{cap:planificacio} i
detalls a la secció corresponent de Disseny).

\paragraph{Recorregut funcional.} La demostració recorre les pantalles
principals del client (vegeu les figures de la secció de disseny):

\begin{enumerate}
    \item Connexió de cartera (Pera Wallet o Lute).
    \item Creació d'una organització i càrrega de cens via \ac{CSV}
    (amb trossejament automàtic en lots de set adreces).
    \item Redacció d'una proposta mitjançant l'assistent de cinc
    passos.
    \item Vot d'aprovació binari per part dels membres del cens.
    \item Vot d'elecció amb \emph{rànquing} arrossegable (component
    \texttt{@dnd-kit}, accessible amb ratolí i teclat).
    \item Visualització dels resultats finals amb el mètode de
    Schulze, la matriu de comparació per parelles col·lapsable i el
    panell de verificació pública del vot.
\end{enumerate}

\paragraph{Verificabilitat.} La demostració inclou la verificació
externa del resultat: introduint una adreça arbitrària al panell de
verificació, el client consulta directament la cadena i mostra el
\emph{rànquing} emès per aquesta adreça, sense necessitat de cap
permís. En paral·lel, el servei d'ancoratge publica el \emph{hash}
determinista de l'elecció tancada al contracte \texttt{NotaryContract}
de Sepolia, fent visible el consens K-de-N entre nodes universitaris.

\paragraph{Temps total.} Sobre el \emph{LocalNet}, el cicle complet
--des de \texttt{docker compose up} fins a la visualització dels
resultats-- es pot executar en menys de \textbf{cinc minuts}, fet que
n'ha facilitat l'ús durant les demostracions de cada \emph{sprint} i
la repetició davant un públic acadèmic a l'entrega final.
